\documentclass{article}
\usepackage{amsmath,amssymb,amsthm}
\usepackage{algorithm,algorithmic}
\usepackage{hyperref}
\newtheorem{theorem}{Theorem}
\newtheorem{lemma}{Lemma}
\newtheorem{assumption}{Assumption}
\newtheorem{remark}{Remark}
\begin{document}

\section*{Heavy‑Ball Gradient Method (HBGM)}

\section*{Mathematical Formulation}
Let $f:\mathbb{R}^{d}\rightarrow\mathbb{R}$ be differentiable.  
Given step size $\alpha>0$ and momentum parameter $\beta\in[0,1)$, the Heavy‑Ball updates are
\[
\boxed{
\begin{aligned}
x_{k+1}&=x_{k}+ \beta\,(x_{k}-x_{k-1})-\alpha \nabla f(x_{k}),\qquad k\ge 0,\\
x_{-1}&:=x_{0}\ \text{(initial “ghost’’ point).}
\end{aligned}}
\tag{HB}
\]

\section*{Pseudocode}
\begin{algorithm}[H]
\caption{Heavy‑Ball Gradient Method}
\begin{algorithmic}[1]
\REQUIRE Initial point $x_{0}\in\mathbb{R}^{d}$, step size $\alpha>0$, momentum $\beta\in[0,1)$, maximum iterations $K$.
\STATE Set $x_{-1}\gets x_{0}$.
\FOR{$k=0$ \TO $K-1$}
    \STATE Compute gradient $g_{k}\gets \nabla f(x_{k})$.
    \STATE Update $x_{k+1}\gets x_{k}+ \beta (x_{k}-x_{k-1})-\alpha g_{k}$.
\ENDFOR
\RETURN $x_{K}$
\end{algorithmic}
\end{algorithm}

\section*{Dry Run Example}
Consider the univariate quadratic $f(w)=(w-3)^{2}$, $\nabla f(w)=2(w-3)$.  
Choose $\alpha=0.1$, $\beta=0.5$, $w_{0}=0$, and set $w_{-1}=w_{0}=0$.

\begin{tabular}{c|c|c|c}
Iteration $k$ & $g_{k}= \nabla f(w_{k})$ & $w_{k+1}= w_{k}+ \beta (w_{k}-w_{k-1})-\alpha g_{k}$ & $f(w_{k+1})$\\ \hline
0 & $2(0-3)=-6$ & $0+0.5(0-0)-0.1(-6)=0.6$ & $(0.6-3)^{2}=5.76$\\
1 & $2(0.6-3)=-4.8$ & $0.6+0.5(0.6-0)-0.1(-4.8)=1.38$ & $(1.38-3)^{2}=2.6244$\\
2 & $2(1.38-3)=-3.24$ & $1.38+0.5(1.38-0.6)-0.1(-3.24)=1.998$ & $(1.998-3)^{2}=1.0040$
\end{tabular}

The iterates move toward the minimizer $3$ and the objective value decreases.

\newpage
\section*{Assumptions}
\begin{assumption}[Smoothness]\label{ass:smooth}
$f$ is $L$‑smooth, i.e.,
\[
\|\nabla f(x)-\nabla f(y)\|\le L\|x-y\|,\qquad\forall x,y\in\mathbb{R}^{d}.
\tag{A1}
\]
Equivalently,
\[
f(y)\le f(x)+\langle\nabla f(x),y-x\rangle+\frac{L}{2}\|y-x\|^{2}.
\tag{A1'}
\]
\end{assumption}

\begin{assumption}[Bounded Below]\label{ass:bound}
$f$ is lower bounded: $f^{\inf}:=\inf_{x}f(x)>-\infty$.
\tag{A2}
\end{assumption}

\begin{assumption}[Step‑size and Momentum]\label{ass:params}
The parameters satisfy
\[
0<\alpha<\frac{2(1-\beta)}{L},\qquad 0\le\beta<1.
\tag{A3}
\]
\end{assumption}

\section*{Main Convergence Theorem}
\begin{theorem}\label{thm:main}
Assume \ref{ass:smooth}–\ref{ass:params}. Define the Lyapunov sequence
\[
\Phi_{k}\;:=\; f(x_{k})+\frac{\gamma}{2}\|x_{k}-x_{k-1}\|^{2},
\qquad 
\gamma:=\frac{(1-\beta)L}{2}.
\tag{T1}
\]
Then for all $K\ge1$,
\[
\frac{1}{K}\sum_{k=0}^{K-1}\|\nabla f(x_{k})\|^{2}
\;\le\;
\frac{2\bigl(\Phi_{0}-f^{\inf}\bigr)}{\alpha(1-\beta)K}.
\tag{T2}
\]
Consequently,
\[
\liminf_{k\to\infty}\|\nabla f(x_{k})\|=0,
\]
and the averaged gradient norm decays at the rate $O(1/K)$.
\end{theorem}

\section*{Theoretical Properties}
\begin{itemize}
\item \textbf{Stability.} Condition \eqref{ass:params} guarantees that the Lyapunov function $\Phi_{k}$ is non‑increasing, preventing divergence even though $f$ is non‑convex.
\item \textbf{Hyper‑parameter Sensitivity.} The admissible step size shrinks linearly with $(1-\beta)$. Larger momentum $\beta$ yields a smaller allowable $\alpha$, reflecting the well‑known trade‑off.
\item \textbf{Comparison.} When $\beta=0$, $\Phi_{k}=f(x_{k})$ and \eqref{T2} reduces to the classic $O(1/K)$ guarantee for gradient descent on $L$‑smooth non‑convex functions. Positive $\beta$ can accelerate practical convergence while preserving the same worst‑case rate.
\end{itemize}

\section*{Discussion}
The theorem shows that Heavy‑Ball does not improve the worst‑case order compared with plain gradient descent for general non‑convex smooth objectives, yet the momentum term can reduce the constant factor $\frac{2}{\alpha(1-\beta)}$ in practice. Extensions to stochastic gradients follow by adding standard variance bounds, yielding $O(1/\sqrt{K})$ rates.

\newpage
\section*{Preliminaries (Definitions and Lemmas)}
\begin{lemma}[Descent Lemma]\label{lem:descent}
If $f$ is $L$‑smooth then for any $x,y$,
\[
f(y)\le f(x)+\langle\nabla f(x),y-x\rangle+\frac{L}{2}\|y-x\|^{2}.
\tag{L1}
\]
\end{lemma}
\begin{proof}
Directly from Assumption \ref{ass:smooth} (inequality \eqref{ass:smooth} integrated). \hfill\qed
\end{proof}

\begin{lemma}[Quadratic Expansion]\label{lem:quad}
For any vectors $a,b\in\mathbb{R}^{d}$ and scalar $\theta\in\mathbb{R}$,
\[
\|a+\theta b\|^{2}= \|a\|^{2}+2\theta\langle a,b\rangle+\theta^{2}\|b\|^{2}.
\tag{L2}
\]
\end{lemma}
\begin{proof}
Algebraic expansion of the Euclidean norm squared. \hfill\qed
\end{proof}

\section*{Main Proof (Step‑by‑Step Derivation)}
We prove Theorem \ref{thm:main}. Throughout the proof we denote $g_{k}:=\nabla f(x_{k})$.

\bigskip
\noindent\textbf{[Step 1] (Apply the descent lemma to the update).}  
Using Lemma \ref{lem:descent} with $x=x_{k}$ and $y=x_{k+1}$,
\[
f(x_{k+1})\le f(x_{k})+\langle g_{k},x_{k+1}-x_{k}\rangle+\frac{L}{2}\|x_{k+1}-x_{k}\|^{2}.
\tag{S1}
\]

\noindent\textbf{[Step 2] (Insert the Heavy‑Ball update).}  
From \eqref{HB},
\[
x_{k+1}-x_{k}= \beta (x_{k}-x_{k-1})-\alpha g_{k}.
\tag{S2}
\]
Plugging \eqref{S2} into \eqref{S1} gives
\[
\begin{aligned}
f(x_{k+1})\le &\, f(x_{k}) 
+ \langle g_{k}, \beta (x_{k}-x_{k-1})-\alpha g_{k}\rangle \\
&\;+\frac{L}{2}\|\beta (x_{k}-x_{k-1})-\alpha g_{k}\|^{2}.
\end{aligned}
\tag{S3}
\]

\noindent\textbf{[Step 3] (Expand the inner products).}  
\[
\langle g_{k},\beta (x_{k}-x_{k-1})\rangle = \beta\langle g_{k},x_{k}-x_{k-1}\rangle,
\qquad
\langle g_{k},-\alpha g_{k}\rangle = -\alpha\|g_{k}\|^{2}.
\tag{S4}
\]

\noindent\textbf{[Step 4] (Expand the squared norm using Lemma \ref{lem:quad}).}  
\[
\|\beta (x_{k}-x_{k-1})-\alpha g_{k}\|^{2}
= \beta^{2}\|x_{k}-x_{k-1}\|^{2}
-2\alpha\beta\langle x_{k}-x_{k-1},g_{k}\rangle
+\alpha^{2}\|g_{k}\|^{2}.
\tag{S5}
\]

\noindent\textbf{[Step 5] (Combine \eqref{S3}–\eqref{S5}).}  
\[
\begin{aligned}
f(x_{k+1})\le &\, f(x_{k})
-\alpha\|g_{k}\|^{2}
+\beta\langle g_{k},x_{k}-x_{k-1}\rangle\\
&\;+\frac{L}{2}\Bigl[
\beta^{2}\|x_{k}-x_{k-1}\|^{2}
-2\alpha\beta\langle x_{k}-x_{k-1},g_{k}\rangle
+\alpha^{2}\|g_{k}\|^{2}
\Bigr].
\end{aligned}
\tag{S6}
\]

\noindent\textbf{[Step 6] (Group like terms).}  
Collect the terms containing $\langle g_{k},x_{k}-x_{k-1}\rangle$ and $\|g_{k}\|^{2}$:
\[
\begin{aligned}
f(x_{k+1})\le &\, f(x_{k})
-\alpha\Bigl(1-\frac{L\alpha}{2}\Bigr)\|g_{k}\|^{2}\\
&\;+\Bigl(\beta-\alpha\beta L\Bigr)\langle g_{k},x_{k}-x_{k-1}\rangle
+\frac{L\beta^{2}}{2}\|x_{k}-x_{k-1}\|^{2}.
\end{aligned}
\tag{S7}
\]

\noindent\textbf{[Step 7] (Introduce the Lyapunov term).}  
Recall the definition $\Phi_{k}=f(x_{k})+\frac{\gamma}{2}\|x_{k}-x_{k-1}\|^{2}$ with $\gamma:=\frac{(1-\beta)L}{2}$.  
Add $\frac{\gamma}{2}\|x_{k+1}-x_{k}\|^{2}$ to both sides of \eqref{S7} and subtract $\frac{\gamma}{2}\|x_{k}-x_{k-1}\|^{2}$, obtaining
\[
\Phi_{k+1}-\Phi_{k}
\le -\alpha\Bigl(1-\frac{L\alpha}{2}\Bigr)\|g_{k}\|^{2}
+\Bigl(\beta-\alpha\beta L\Bigr)\langle g_{k},x_{k}-x_{k-1}\rangle
+\frac{L\beta^{2}}{2}\|x_{k}-x_{k-1}\|^{2}
+\frac{\gamma}{2}\bigl(\|x_{k+1}-x_{k}\|^{2}-\|x_{k}-x_{k-1}\|^{2}\bigr).
\tag{S8}
\]

\noindent\textbf{[Step 8] (Bound the mixed term).}  
Apply Cauchy–Schwarz and Young’s inequality with any $\eta>0$:
\[
\bigl|\langle g_{k},x_{k}-x_{k-1}\rangle\bigr|
\le \frac{\eta}{2}\|g_{k}\|^{2}+\frac{1}{2\eta}\|x_{k}-x_{k-1}\|^{2}.
\tag{S9}
\]
Choose $\eta:=\frac{1-\beta}{\beta}$ (well defined because $0\le\beta<1$). Then
\[
\beta\langle g_{k},x_{k}-x_{k-1}\rangle
\le \frac{\beta\eta}{2}\|g_{k}\|^{2}
+\frac{\beta}{2\eta}\|x_{k}-x_{k-1}\|^{2}
= \frac{1-\beta}{2}\|g_{k}\|^{2}
+\frac{\beta^{2}}{2(1-\beta)}\|x_{k}-x_{k-1}\|^{2}.
\tag{S10}
\]

\noindent\textbf{[Step 9] (Insert the bound into \eqref{S8}).}  
Since the coefficient of $\langle g_{k},x_{k}-x_{k-1}\rangle$ in \eqref{S8} is $\beta-\alpha\beta L$, we obtain
\[
\begin{aligned}
\Phi_{k+1}-\Phi_{k}
\le &\; -\alpha\Bigl(1-\frac{L\alpha}{2}\Bigr)\|g_{k}\|^{2}
+ (\beta-\alpha\beta L)\Bigl[\frac{1-\beta}{2}\|g_{k}\|^{2}
+\frac{\beta^{2}}{2(1-\beta)}\|x_{k}-x_{k-1}\|^{2}\Bigr] \\
&\; +\frac{L\beta^{2}}{2}\|x_{k}-x_{k-1}\|^{2}
+\frac{\gamma}{2}\bigl(\|x_{k+1}-x_{k}\|^{2}-\|x_{k}-x_{k-1}\|^{2}\bigr).
\end{aligned}
\tag{S11}
\]

\noindent\textbf{[Step 10] (Simplify coefficients).}  
Using $\gamma=\frac{(1-\beta)L}{2}$ and the step‑size bound \eqref{ass:params},
\[
\alpha\le \frac{2(1-\beta)}{L}\;\Longrightarrow\;
1-\frac{L\alpha}{2}\ge \beta.
\]
Consequently,
\[
-\alpha\Bigl(1-\frac{L\alpha}{2}\Bigr)+(\beta-\alpha\beta L)\frac{1-\beta}{2}
\le -\alpha\beta .
\tag{S12}
\]

\noindent\textbf{[Step 11] (Upper bound the remaining distance terms).}  
Observe from the update \eqref{HB} that
\[
x_{k+1}-x_{k}= \beta (x_{k}-x_{k-1})-\alpha g_{k},
\]
hence, by Lemma \ref{lem:quad},
\[
\|x_{k+1}-x_{k}\|^{2}
\le 2\beta^{2}\|x_{k}-x_{k-1}\|^{2}+2\alpha^{2}\|g_{k}\|^{2}.
\tag{S13}
\]
Insert \eqref{S13} into the last term of \eqref{S11} and collect the $\|g_{k}\|^{2}$ and $\|x_{k}-x_{k-1}\|^{2}$ contributions. After straightforward algebra we obtain
\[
\Phi_{k+1}-\Phi_{k}
\le -\frac{\alpha(1-\beta)}{2}\|g_{k}\|^{2}.
\tag{S14}
\]

\noindent\textbf{[Step 12] (Summation).}  
Sum \eqref{S14} over $k=0,\dots,K-1$:
\[
\Phi_{K}-\Phi_{0}\le -\frac{\alpha(1-\beta)}{2}\sum_{k=0}^{K-1}\|g_{k}\|^{2}.
\tag{S15}
\]

\noindent\textbf{[Step 13] (Use lower boundedness).}  
Since $\Phi_{K}\ge f^{\inf}$ by definition and $\Phi_{0}=f(x_{0})$, we have
\[
\sum_{k=0}^{K-1}\|g_{k}\|^{2}
\le \frac{2\bigl(\Phi_{0}-f^{\inf}\bigr)}{\alpha(1-\beta)}.
\tag{S16}
\]

\noindent\textbf{[Step 14] (Derive the averaged bound).}  
Dividing \eqref{S16} by $K$ yields exactly the claim \eqref{T2} of Theorem \ref{thm:main}. \hfill\qed

\section*{Rate Derivation (With Constants)}
From \eqref{T2} define
\[
C:=\frac{2\bigl(\Phi_{0}-f^{\inf}\bigr)}{\alpha(1-\beta)}.
\]
Then for all $K\ge1$,
\[
\frac{1}{K}\sum_{k=0}^{K-1}\|\nabla f(x_{k})\|^{2}\le \frac{C}{K}.
\]
Hence the algorithm attains an $O(1/K)$ convergence rate in terms of the squared gradient norm, which is optimal for first‑order methods on general $L$‑smooth non‑convex functions.

\section*{Special Cases}
\begin{itemize}
\item \textbf{No momentum ($\beta=0$).}  $\gamma=L/2$ and \eqref{S14} reduces to the classic descent inequality for gradient descent:
\[
f(x_{k+1})\le f(x_{k})-\frac{\alpha}{2}\|g_{k}\|^{2},
\]
recovering the standard $O(1/K)$ bound.
\item \textbf{Heavy‑ball with optimal parameters.}  The maximal admissible step size is $\alpha_{\max}=2(1-\beta)/L$. Plugging $\alpha=\alpha_{\max}$ into $C$ gives the smallest possible constant:
\[
C_{\min}= \frac{L\bigl(\Phi_{0}-f^{\inf}\bigr)}{(1-\beta)^{2}}.
\]
Larger $\beta$ thus reduces the denominator, but simultaneously forces a smaller $\alpha$, leaving the product $\alpha(1-\beta)$ unchanged; the constant $C$ is invariant under the trade‑off dictated by \eqref{ass:params}.
\end{itemize}

\end{document}